% =====================================================================
% Lampiran — TIDAK dihitung dalam batas 12 halaman.
% Berisi detail yang dipindahkan dari isi utama agar isi utama tetap ketat.
% =====================================================================

\appendix

\section{Konfigurasi Pelatihan dan \emph{Trade-off} Rancangan}
\label{app:konfig}

\paragraph{Praproses.} Masukan diseragamkan ke $224 \times 224$; RGB dinormalkan
dengan statistik ImageNet dan termal dengan rerata/simpangan $0{,}5$. Augmentasi
geometris (pencerminan horizontal, rotasi $\pm 10^\circ$) diterapkan identik pada
RGB, termal, dan mask agar keselarasan spasial terjaga.

\paragraph{Konfigurasi pelatihan.} Pengoptimal AdamW; laju pembelajaran
$3 \times 10^{-4}$; \emph{weight decay} $0{,}01$; penjadwalan \emph{cosine};
ukuran \emph{batch} 16; maksimum 40 epoch dengan penghentian dini berkesabaran
10; benih acak 42. \emph{Modality dropout} $p = 0{,}2$, dipilih dari grid
$\{0{,}1;\ 0{,}2;\ 0{,}3\}$. \emph{Gate} keandalan dilatih 30 epoch dengan bobot
\emph{loss} regresi 1{,}0 terhadap \emph{loss} klasifikasi. \emph{Head}
segmentasi dilatih 30 epoch di atas \emph{encoder} beku, laju pembelajaran
$5 \times 10^{-4}$. Seluruh pelatihan berjalan pada satu sesi GPU kelas
menengah; satu run memakan 20--60 menit.

\paragraph{Penyelarasan grid token.} \emph{Encoder} RGB menghasilkan grid
$16 \times 16$ (256 token) dan \emph{encoder} termal $14 \times 14$ (196 token).
Keduanya diselaraskan lewat interpolasi bilinear sebelum fusi agar indeks token
merujuk lokasi fisik yang sama.

\paragraph{\emph{Trade-off} yang disengaja.}
\begin{itemize}
\item \textbf{Ubin $224 \times 224$ alih-alih $336 \times 336$.} Resolusi lebih
tinggi menaikkan akurasi objek kecil, persis target deteksi dini, tetapi
memperlambat inferensi. Kami memilih kelayakan inferensi.
\item \textbf{Membekukan 9 dari 12 blok \emph{backbone}.} Konvergensi cepat dan
hemat memori pada data terbatas, dengan biaya kapasitas adaptasi domain yang
lebih rendah, justru dimensi yang paling dibutuhkan untuk pindah ke gambut.
\item \textbf{Grid $p$ kecil.} Kompromi terhadap keterbatasan waktu, bukan
pencarian hiperparameter ekstensif.
\item \textbf{Dua lapis fusi.} Demi latensi; fusi tetap menjadi penyumbang
parameter terlatih terbesar (7,1 juta dari 13,9 juta).
\item \textbf{Satu benih acak.} Menukar ketelitian statistik dengan cakupan
eksperimen dalam jendela 12 hari.
\end{itemize}

\section{Definisi Metrik dan Uji Statistik}
\label{app:metrik}

Presisi $P = \mathrm{TP}/(\mathrm{TP}+\mathrm{FP})$, \emph{recall}
$R = \mathrm{TP}/(\mathrm{TP}+\mathrm{FN})$, $\mathrm{F1} = 2PR/(P+R)$,
dilaporkan per kelas beserta rerata makro. Kelas tanpa dukungan dikeluarkan dari
rerata makro; melibatkannya menghasilkan angka yang menyesatkan (lihat
Subbagian~4.2). IoU $= |A \cap B|/|A \cup B|$, diagregasi menjadi mIoU sebagai
rerata antar kelas dengan sampel tanpa kemunculan kelas diabaikan, bukan
dihitung nol. \emph{Pointing game accuracy} mengukur proporsi sampel yang
titik atensi tertingginya jatuh di dalam wilayah anotasi. Selisih ketahanan
$\Delta_m = M(\text{fusi}) - M(\text{modalitas } m)$.

Selang kepercayaan memakai metode Wilson pada tingkat 95\%. Perbandingan
antar-konfigurasi memakai uji McNemar eksak dua sisi atas pasangan diskordan,
karena seluruh konfigurasi dievaluasi pada sampel yang sama; membandingkan
selang kepercayaan yang tumpang tindih akan terlalu konservatif untuk desain
berpasangan. Latensi dilaporkan sebagai median dari sekurang-kurangnya 100
pengukuran setelah 10 kali pemanasan.

\section{Rincian Efisiensi dan Kuantisasi}
\label{app:efisiensi}

\begin{table}[h]
\caption{Ukuran, kompleksitas, dan latensi. Seluruh latensi diukur pada laptop
Lenovo ThinkPad T480, Windows, ONNX Runtime, presisi fp32 kecuali disebutkan
lain, satu bingkai penuh $224 \times 224$.}
\centering
\small
\begin{tabular}{lr}
\hline
\textbf{Besaran} & \textbf{Nilai} \\
\hline
Parameter total & 30.626.467 \\
Parameter terlatih & 13.896.355 \\
FLOPs per \emph{forward} (thop) & 18{,}68 G \\
Ukuran \emph{checkpoint} fp32 & 116{,}94 MB \\
Ukuran ONNX fp32 $\rightarrow$ INT8 & 117{,}10 $\rightarrow$ 36{,}53 MB ($-68{,}8\%$) \\
Latensi ONNX fp32 & 117{,}5 ms median; p95 159{,}3; rentang antar-pemanggilan 101{,}7--140{,}2 \\
Latensi ONNX INT8 & 390{,}8 ms median (\emph{speedup} $0{,}36\times$) \\
Latensi PyTorch \emph{eager}, Intel Core Ultra 7 255H & 44{,}7 ms median; p90 47{,}1 \\
\hline
\end{tabular}
\end{table}

Model beroperasi pada bingkai penuh, bukan per ubin, sehingga angka di atas
sudah merupakan latensi per bingkai dan tidak perlu dikalikan jumlah ubin.
Tidak ada perangkat tepi fisik yang diuji; angka ini adalah proksi CPU x86.
Dampak kuantisasi terhadap akurasi tidak dapat ditetapkan karena evaluasi yang
tersedia dijalankan pada partisi yang dibangun dari rentang bingkai berstatus
sementara dan menghasilkan metrik degeneratif.

\section{Evaluasi Heuristik dan Audit Aksesibilitas}
\label{app:heuristik}

Antarmuka dikembangkan melalui 19 putaran evaluasi terdokumentasi yang
menghasilkan sepuluh temuan, seluruhnya menyebabkan perubahan antarmuka. Empat
temuan hanya dapat ditemukan lewat pengukuran, bukan penilaian visual:

\begin{enumerate}
\item Tiga dari enam warna tingkat risiko gagal ambang kontras WCAG~AA 4{,}5:1
terhadap permukaan kartu (2{,}93:1; 3{,}24:1; 3{,}53:1). Perbaikan memisahkan
tiap tingkat risiko menjadi dua nilai, satu untuk isian, satu khusus teks.
\item Simulasi buta warna dengan perhitungan $\Delta E$ atas seluruh pasangan
seri menemukan kombinasi hijau-tua dan merah-tua menghasilkan $\Delta E = 4{,}3$
pada protanopia. Salah satu seri diganti, dan identitas seri tidak lagi
bergantung warna semata.
\item Median sigmoid keluaran \emph{head} segmentasi berada di sekitar $0{,}69$,
sehingga tumpang-tindih lokalisasi menutupi hampir seluruh bingkai. Ambang
persentil-90 diterapkan sebelum render.
\item Selisih keyakinan pada pita kartu tidak terbaca pada rentang tebal
4--10~px; rentang dinaikkan dan ditambahkan takik pada batas isian.
\end{enumerate}

Satu cacat diperkenalkan oleh evaluator sendiri pada putaran kedua dan
tertangkap pada putaran ketiga, dicatat apa adanya karena menunjukkan putaran
evaluasi berulang memang menangkap regresi, bukan sekadar memoles.

\section{Estimasi Biaya Penerapan}
\label{app:biaya}

\begin{table}[h]
\caption{Estimasi biaya penerapan. Harga ritel merupakan daftar publik yang
dapat berubah dan perlu diverifikasi ulang sebelum dipakai sebagai dasar
pengadaan.}
\centering
\small
\begin{tabular}{lll}
\hline
\textbf{Komponen} & \textbf{Estimasi} & \textbf{Sifat} \\
\hline
Drone RGB (armada eksisting) & Rp 0 tambahan & modal, sudah ada \\
Drone RGB+termal, mis.\ DJI Mavic 3T & Rp 79{,}55 juta/unit & modal, opsional \\
Perangkat tepi kelas Jetson Orin Nano & Rp 9{,}5--12{,}9 juta/unit & modal, sekali beli \\
Pelatihan model, satu siklus & Rp 0 (kuota GPU gratis) & operasional \\
Inferensi operasional & Rp 0 marginal (\emph{on-device}) & operasional \\
API hotspot satelit & Rp 0 & operasional \\
Hosting papan pantau & Rp 0--156 ribu/bulan & berulang \\
\hline
\end{tabular}
\end{table}

Harga drone termal merupakan daftar ritel resmi Indonesia per Agustus 2026
(Galaxy Camera, mitra resmi DJI, sudah termasuk PPN); harga perangkat tepi
merupakan rentang daftar ritel pada Blibli dan Tokopedia. Keduanya dapat berubah
dan wajib diverifikasi ulang sebelum dipakai sebagai dasar pengadaan.

Belum dihitung dan dinyatakan terbuka: biaya operator dan pelatihan personel,
perawatan serta penggantian baterai drone, perizinan penerbangan, dan integrasi
ke sistem instansi yang sudah berjalan.

\section{Integrasi Pemicu Satelit}
\label{app:firms}

Lapisan pemicu menarik data hotspot NASA FIRMS VIIRS SNPP NRT melalui proksi
sisi server pada bounding box Sumatra tengah--selatan dan Kalimantan
($-4{,}5^\circ$ s.d.\ $+1{,}5^\circ$ lintang; $100^\circ$ s.d.\ $117^\circ$
bujur). VIIRS dipilih ketimbang MODIS karena resolusi spasialnya 375~m
dibanding 1~km, sehingga titik panas kecil yang khas pada kebakaran gambut
tahap awal berpeluang lebih besar terdeteksi.

Respons dipangkas ke 400 titik ber-FRP tertinggi karena pada puncak musim
kemarau area tersebut mengembalikan lebih dari 2.000 titik yang menenggelamkan
\emph{alert} model di peta; jumlah total tetap dinyatakan di antarmuka. Pada
verifikasi 9~Agustus~2026, 2.467 titik terdeteksi dengan FRP 15{,}1--128{,}9~MW,
waktu respons 1{,}03 detik tanpa \emph{cache} dan 0{,}14 detik dengan
\emph{cache}. Kunci API disimpan sebagai variabel lingkungan di sisi server,
tidak pernah dikirim ke peramban maupun masuk log.

Keterbatasan: data bersifat \emph{near real-time} (tersedia beberapa jam setelah
lintasan), resolusi 375~m dapat melewatkan kebakaran bawah permukaan, VIIRS
bersifat optik sehingga terhalang awan tebal, hotspot menandai anomali termal
dan bukan kebakaran, belum ada validasi silang antara hotspot dan keluaran
model, area pantau bersifat tetap, dan kuota berlaku per kunci API.

\section{Antarmuka Purwarupa}
\label{app:antarmuka}

Konsol operator terdiri atas empat halaman: peta operasi (hotspot satelit,
penanda \emph{alert}, sensor simulasi berlabel, rute patroli), panel
\emph{alert} (daftar terurut risiko dengan lencana modalitas di tiap kartu),
rincian \emph{alert} (RGB dan termal berdampingan, tumpang-tindih lokalisasi,
keandalan modalitas, jejak keputusan), dan halaman informasi model untuk
pembaca teknis (tabel hasil, kurva degradasi, metrik lokalisasi).

Tangkapan layar diambil dari \emph{build} produksi dengan data nyata: seluruh
\emph{alert} merupakan keluaran model, hotspot berasal dari layanan satelit
langsung, dan angka evaluasi berasal dari berkas hasil eksperimen. Catatan
provenans pada halaman rincian menyatakan bahwa bingkai citra berasal dari
FLAME~2 (Arizona) dan koordinat merupakan penempatan simulasi.

\section{Risiko dan Mitigasi}
\label{app:risiko}

\begin{table}[!htbp]
\caption{Risiko dan mitigasi. Status membedakan yang sudah diterapkan pada
purwarupa dari yang masih berupa rancangan.}
\label{tab:risiko}
\centering
\small
\begin{tabular}{p{1.7cm}p{4.6cm}p{4.4cm}p{1.4cm}}
\hline
\textbf{Risiko} & \textbf{Konteks} & \textbf{Mitigasi} & \textbf{Status} \\
\hline
Privasi & Patroli melintasi permukiman dan lahan garapan masyarakat &
Pembatasan masa simpan citra; pengaburan wajah dan pelat kendaraan; jalur
patroli menghindari permukiman & rancangan \\
Bias & Dilatih pada hutan pinus beriklim sedang; kelas api-tanpa-asap tidak
pernah terpelajari &
Batas domain dinyatakan terbuka di paper dan antarmuka; adaptasi domain ringan
per wilayah & sebagian \\
Halusinasi & Tidak berlaku langsung, model diskriminatif, bukan generatif.
Risiko analognya keyakinan berlebih di luar distribusi &
Skor keandalan per modalitas ditampilkan ke operator; batas domain eksplisit;
kalibrasi formal belum & sebagian \\
Keamanan & Lokasi kebakaran di areal konsesi sensitif secara komersial dan
hukum; API dan papan pantau adalah permukaan serangan &
Kunci API di sisi server, tidak pernah masuk log; galat tidak bocor ke layar
operator. Autentikasi, kontrol akses berbasis peran, dan log akses belum ada &
sebagian \\
Keselamatan & Kegagalan mendeteksi jauh lebih mahal daripada alarm palsu &
Titik operasi condong ke \emph{recall}; fusi menekan alarm palsu termal ke nol;
keputusan operator jadi umpan balik kalibrasi & sudah \\
Penyalah\-gunaan & Dapat dialihfungsikan jadi alat pengawasan terhadap
masyarakat sekitar &
Keluaran dibatasi pada deteksi kebakaran; log akses; pemisahan kendali data
antar pihak & rancangan \\
\hline
\end{tabular}
\end{table}

% =====================================================================
% Lampiran tambahan — menutup sub-kriteria rubrik yang belum punya
% tempat eksplisit: 1a (pengguna & pain point), 2a/2b/2c (adopsi,
% kelayakan, skalabilitas), 4a (alur pengguna ujung-ke-ujung),
% serta diagram sistem dari concept paper penyisihan.
%
% Disisipkan di lampiran.tex, TIDAK dihitung dalam batas 12 halaman.
% =====================================================================

\section{Pengguna Sasaran dan Titik Nyeri Operasional}
\label{app:pengguna}

\subsection*{Segmen pengguna}

\begin{table}[!htbp]
\centering
\small
\begin{tabular}{p{2.6cm}p{4.2cm}p{5.2cm}}
\hline
\textbf{Segmen} & \textbf{Siapa} & \textbf{Yang berubah bila sistem dipakai} \\
\hline
Instansi penanggulangan &
Manggala Agni, BNPB, dinas kehutanan daerah &
Dasar objektif untuk memilih hotspot mana yang diverifikasi lebih dulu, sebelum
tim dan bahan bakar dikeluarkan \\
Pemegang konsesi &
Perusahaan perkebunan sawit dan hutan tanaman industri di lanskap gambut &
Bukti pemantauan mandiri atas areal yang menjadi tanggung jawab mutlaknya
menurut Pasal~88 UU~32/2009 \\
Operator papan pantau &
Petugas posko, bukan tenaga pembelajaran mesin &
Antarmuka berbahasa Indonesia dengan prioritas terurut, keandalan per sensor,
dan tiga tombol keputusan \\
\hline
\end{tabular}
\caption{Segmen pengguna sasaran dan perubahan yang dijanjikan sistem.}
\end{table}

\subsection*{Titik nyeri, sumbernya, dan jawaban sistem}

Empat titik nyeri di bawah berasal dari wawancara narasumber ahli
(Bagian~5) dan dari karakteristik infrastruktur pemantauan yang berlaku umum.
Titik nyeri yang tidak memiliki sumber terverifikasi tidak kami cantumkan.

\begin{enumerate}
\item \textbf{Hotspot satelit memicu, tetapi tidak mengonfirmasi.} FIRMS
menandai anomali termal beresolusi 375~m; tungku industri, suar gas, dan tanah
panas ikut terdeteksi. Tim lapangan tidak punya dasar objektif untuk memilih
mana yang dikejar duluan. \emph{Jawaban sistem}: lapisan konfirmasi visual
beresolusi tinggi yang mengubah anomali termal menjadi penilaian ber-keyakinan,
bukan sekadar titik di peta.

\item \textbf{Drone dipakai setelah kebakaran diketahui, bukan sebelumnya.}
Narasumber menyatakan pemanfaatan drone selama ini terbatas pada pengambilan
citra berkala, misalnya tiap satu jam, sesudah kebakaran terjadi, dengan
penyebaran api diprediksi manual. \emph{Jawaban sistem}: menempatkan model pada
fase verifikasi awal, sehingga keluarannya menjadi masukan keputusan pengerahan,
bukan dokumentasi pasca-kejadian.

\item \textbf{Daya tahan baterai membatasi cakupan patroli.} Sekitar 30 menit
tanpa beban tambahan, dan lebih singkat bila drone juga mengirim data mentah.
\emph{Jawaban sistem}: inferensi berjalan di perangkat dan yang dikirim hanya
\emph{alert}, bukan aliran video, sehingga anggaran daya untuk komunikasi
ditekan.

\item \textbf{Sensor dapat gagal justru saat paling dibutuhkan.} Asap pekat
membutakan kanal optik; \emph{drift} kalibrasi mengganggu kanal termal; sebagian
armada tidak membawa sensor termal sama sekali. \emph{Jawaban sistem}: pelatihan
dengan \emph{modality dropout} sehingga sistem tetap memberi keluaran saat satu
kanal hilang, disertai skor keandalan yang menyatakan sedang bersandar pada
sensor mana.
\end{enumerate}

\section{Jalur Adopsi, Kelayakan Implementasi, dan Skalabilitas}
\label{app:adopsi}

\subsection*{Jalur adopsi, empat tahap}

Jalur di bawah disusun agar tiap tahap menghasilkan sesuatu yang dapat dinilai
sebelum tahap berikutnya dibiayai, bukan sebagai rencana sekali jalan.

\begin{enumerate}
\item \textbf{Menumpang alur yang sudah berjalan.} Instansi sudah memakai
FIRMS dan SiPongi. Sistem ini tidak menggantikannya, melainkan menambahkan
lapisan konfirmasi sesudah pemicu. Tidak ada proses eksisting yang harus
dihentikan agar sistem dapat dipakai.
\item \textbf{Pilot satu posko.} Satu posko, armada drone yang sudah dimiliki,
papan pantau di peramban. Keluaran yang dinilai: berapa banyak hotspot yang
tersaring sebelum tim berangkat, dan seberapa sering operator menyetujui
penilaian sistem.
\item \textbf{Perluasan wilayah.} Area pantau bersifat parameter, bukan kode,
sehingga penambahan wilayah tidak memerlukan perubahan model.
\item \textbf{Adaptasi lanskap.} Lanskap dengan karakteristik visual berbeda
ditangani melalui siklus adaptasi domain ringan per wilayah, bukan pelatihan
ulang dari nol.
\end{enumerate}

Peran operasional dipisahkan tegas: operator posko menjalankan sistem harian dan
mengambil keputusan; tenaga teknis hanya dibutuhkan pada tahap pemasangan dan
siklus adaptasi. Sistem tidak menuntut instansi memelihara kompetensi
pembelajaran mesin secara permanen.

\subsection*{Kelayakan implementasi}

\begin{table}[!htbp]
\centering
\small
\begin{tabular}{p{3.0cm}p{8.8cm}}
\hline
\textbf{Aspek} & \textbf{Status terverifikasi} \\
\hline
Anggaran komputasi &
Seluruh pelatihan Stage 1--7 berjalan pada kuota GPU gratis, satu run 20--60
menit. Arsitektur memang diturunkan dari asumsi satu GPU kelas menengah:
\emph{backbone} kecil, 9 dari 12 blok dibekukan, ubin $224 \times 224$. \\
Ketersediaan data &
Seluruh data pelatihan berasal dari sumber publik gratis dengan lisensi jelas
(FLAME~2 CC~BY~4.0, anotasi RFFNet MIT, FLAME~3). Tidak ada dependensi data
berbayar. \\
Kebutuhan inferensi &
117,5~ms per bingkai pada CPU laptop kelas menengah tanpa akselerator, presisi
fp32. Berjalan tanpa GPU. \\
Dependensi eksternal &
Hanya API hotspot satelit publik gratis dan basemap peta. Kegagalan salah
satunya tidak melumpuhkan halaman: lapisan lain tetap tersaji. \\
Kebutuhan jaringan &
Inferensi tidak memerlukan jaringan. Jaringan hanya dibutuhkan untuk lapisan
hotspot dan basemap. \\
Kendala hukum dan privasi &
Dibahas pada Bagian~6; mitigasi yang belum berjalan dinyatakan sebagai
rancangan, bukan disajikan sebagai perlindungan yang sudah ada. \\
\hline
\end{tabular}
\caption{Aspek kelayakan implementasi dan status verifikasinya.}
\end{table}

\subsection*{Skalabilitas}

Skalabilitas kami nyatakan berdasarkan tiga hasil yang terukur, bukan janji.

\textbf{Pertama, biaya marginal per patroli mendekati nol.} Inferensi berjalan
di perangkat sendiri tanpa API berbayar, sehingga menambah frekuensi atau
cakupan patroli tidak menambah biaya perangkat lunak sama sekali. Biaya
berulang satu-satunya adalah hosting papan pantau.

\textbf{Kedua, perpindahan ke dataset dan lokasi lain tidak menuntut pelatihan
ulang.} Akurasi \emph{zero-shot} 92\% ke FLAME~3, tanpa satu pun sampel
adaptasi, menunjukkan representasi yang dipelajari tidak terkunci pada satu
himpunan data. Batasnya jelas dan kami nyatakan: kedua dataset berada di
Amerika Serikat dan beriklim sedang, sehingga angka ini bukan bukti transfer ke
gambut tropis.

\textbf{Ketiga, perluasan ke lanskap baru tidak menuntut anotasi piksel.} Jalur
lokalisasi tanpa label mencapai \emph{pointing game accuracy} 0,6133 pada
partisi uji. Anotasi piksel adalah komponen termahal dalam menyiapkan wilayah
baru, dan hasil ini menunjukkan lokalisasi tetap berjalan tanpanya, meski dengan
mutu yang lebih rendah daripada jalur bersupervisi penuh.

Yang membatasi skalabilitas dan tidak kami tutupi: adaptasi \emph{few-shot}
dengan LoRA menghasilkan kurva yang tidak monoton dan tidak stabil antar-run,
sehingga jumlah sampel yang dibutuhkan untuk pindah ke domain baru
\textbf{belum dapat kami tetapkan}.

\section{Alur Pengguna Ujung-ke-Ujung}
\label{app:alur}

Alur inti berikut berjalan penuh pada purwarupa dan dapat ditelusuri pada
\emph{deployment}.

\begin{enumerate}
\item \textbf{Pemicu.} Hotspot satelit ditarik langsung untuk area pantau
Sumatra dan Kalimantan, atau penugasan patroli terjadwal. Titik ditampilkan di
peta operasi dengan pembeda sumber dan legenda tetap.
\item \textbf{Penilaian.} Bingkai RGB dan termal dari patroli diproses model
fusi. Keluarannya mengikuti satu kontrak tunggal: label dan keyakinan prediksi,
skor keandalan per modalitas, jalur peta lokalisasi beserta metodenya, sumber
pemicu, dan medan keputusan operator.
\item \textbf{Prioritas.} Panel \emph{alert} menampilkan daftar terurut menurut
risiko. Tiap kartu membawa lencana modalitas, sehingga operator langsung tahu
penilaian itu bersandar pada sensor mana.
\item \textbf{Verifikasi.} Halaman rincian menyandingkan bingkai RGB dan termal
dengan tumpang-tindih lokalisasi yang dapat diatur opasitasnya, keandalan per
modalitas, dan catatan provenans citra.
\item \textbf{Keputusan.} Operator memilih satu dari tiga tindakan:
ditindaklanjuti, ditunda, atau alarm palsu. Pilihan tercatat pada jejak
keputusan dan menjadi umpan balik untuk kalibrasi ambang.
\end{enumerate}

\paragraph{Penanganan kegagalan.} Empat kondisi gagal ditangani dan dapat
diperagakan: layanan hotspot tidak merespons menghasilkan pesan berbahasa
Indonesia disertai penyajian salinan terakhir yang berhasil, dan peta tetap
dapat dipakai karena lapisan lain tidak terpengaruh; galat parsing menghasilkan
pesan fungsional, bukan jejak galat teknis di layar operator; citra yang tidak
tersedia menghasilkan penampung, bukan gambar rusak; dan hilangnya satu
modalitas \textbf{bukan kondisi galat} melainkan jalur yang dirancang, sistem
tetap memberi prediksi dan lencana keandalan berubah menyatakan pergeseran
sandaran.

Butir terakhir adalah titik temu antara kontribusi teknis dan pengalaman
pengguna: perilaku yang pada sistem lain akan menjadi kegagalan, di sini menjadi
keluaran yang dapat dibaca operator.

\section{Rancangan Awal dan Perubahannya setelah Eksperimen}
\label{app:diagram}

Kedua diagram pada lampiran ini adalah \textbf{rancangan yang diajukan pada
\emph{concept paper} penyisihan}, disertakan tanpa penyuntingan agar rancangan
awal dan hasil akhir dapat dibandingkan langsung. Keduanya
\textbf{bukan gambaran sistem sebagaimana terwujud}. Tabel~\ref{tab:perubahan}
memerinci setiap komponen yang berubah, tidak terlaksana, atau bertambah, dan
tabel itulah yang berlaku bila terjadi perbedaan dengan gambar.

\begin{table}[!htbp]
\centering
\small
\begin{tabular}{p{4.6cm}p{7.6cm}}
\hline
\textbf{Pada rancangan awal} & \textbf{Yang benar-benar terjadi} \\
\hline
\multicolumn{2}{l}{\emph{Data dan protokol}} \\
Dataset FLAME, FLAME~2, dan D-Fire &
FLAME~2 untuk pra-pelatihan, anotasi piksel RFFNet untuk \emph{fine-tuning},
FLAME~3 untuk uji kesenjangan domain. FLAME~1 dan D-Fire tidak dipakai \\
Pemisahan per-adegan &
Partisi resmi RFFNet berbasis daftar, dengan anti-kebocoran lewat eksklusi 440
indeks bingkai dari himpunan latih FLAME~2 \\
\hline
\multicolumn{2}{l}{\emph{Model}} \\
\emph{Head} deteksi \emph{anchor-free} dengan keluaran kotak pembatas &
Tidak dibangun. Keluaran model adalah label tiga kelas beserta keyakinan, dan
masker segmentasi \\
Pembanding \emph{backbone} ViT-B/16 (Ablation~\#4) &
Tidak dijalankan; Ablation~\#4 tidak terlaksana \\
\emph{Encoder} termal $\pm 2{,}5$ juta parameter &
1.273.256 parameter \\
\emph{Encoder} RGB $\pm 21$ juta ($\pm 5$ juta terlatih) &
22.056.576 parameter, 5.326.464 di antaranya terlatih \\
Tidak tergambar &
\textbf{Bertambah}: \emph{reliability gate} (Stage~5, 115.202 parameter), dua
jalur lokalisasi (Stage~6), dan adaptasi domain LoRA (Stage~7) \\
\hline
\multicolumn{2}{l}{\emph{Masukan dan penerapan}} \\
Kotak "informasi lain (cuaca, riwayat lahan)" pada lapisan verifikasi &
\textbf{Tidak diimplementasikan.} Model hanya menerima citra RGB dan termal;
variabel lingkungan bukan masukan, sesuai batasan~(ii) pada
Subbagian~\ref{sec:batasan} \\
Masukan berupa ubin $224 \times 224$ &
Model beroperasi pada bingkai penuh, sehingga latensi per-ubin sama dengan
latensi per-bingkai \\
Target latensi $<100$ ms/ubin pada perangkat tepi kelas Jetson, diverifikasi
melalui Ablation~\#4 &
\textbf{Tidak ada perangkat Jetson yang diuji} dan Ablation~\#4 tidak
dijalankan. Yang terukur: 117{,}5 ms per bingkai, ONNX fp32, laptop Lenovo
ThinkPad T480 \\
Mesin data sintetis \emph{diffusion} (ditandai \emph{roadmap}) &
Tetap di luar jalur kritis dan tidak dikerjakan, sesuai rencana \\
\hline
\end{tabular}
\caption{Perbedaan antara rancangan pada \emph{concept paper} penyisihan dan
sistem yang terwujud. Bila gambar dan tabel ini berbeda, tabel ini yang
berlaku.}
\label{tab:perubahan}
\end{table}

Kami menyajikan perbandingan ini secara terbuka alih-alih menyunting gambar
lama, karena selisih antara rancangan dan hasil adalah bagian dari laporan:
empat komponen tidak terlaksana, tiga komponen bertambah, dan satu asumsi
masukan dibatalkan.

\begin{figure}[!htbp]
\centering
\includegraphics[width=\textwidth]{diagramalureksperimen.png}
\caption{	extbf{Rancangan awal} alur eksperimen sebagaimana diajukan pada
\emph{concept paper} penyisihan, bukan alur sebagaimana terlaksana. Komponen
yang tidak terlaksana (pembanding ViT-B/16, Ablation~\#4, \emph{head} deteksi)
dan yang bertambah dirinci pada Tabel~
ef{tab:perubahan}.}
\label{fig:alur}
\end{figure}

\begin{figure}[!htbp]
\centering
\includegraphics[width=\textwidth]{arsitektursistem.png}
\caption{	extbf{Rancangan awal} arsitektur sistem dan \emph{pipeline} model
dari \emph{concept paper} penyisihan. Bagian atas menunjukkan alur keputusan
operasional dari pemicu satelit hingga penerusan \emph{alert}; panel bawah
memerinci jalur model. Jumlah parameter, \emph{head} deteksi, kotak masukan
cuaca, dan target latensi Jetson pada gambar ini 	extbf{tidak terwujud atau
berbeda}; lihat Tabel~
ef{tab:perubahan}.}
\label{fig:arsitektur}
\end{figure}

